\section{Conclusion\label{sec:conclusion}}
This selective survey has drawn connections between
the literature on structural estimation of dynamic discrete choice (DDC) models and the literature on inverse reinforcement learning (IRL), 
and their common ancestor:  Markovian Decision Processes (MDP).
MDPs provide a recipe for how to model the behavior of an intelligent agent who makes sequential
decisions to maximize the expected, discounted value of a stream
of rewards.  Though MDPs can be viewed as specialized in some respects and some of their
assumptions (expected utility, discounted utility) are inconsistent with empirical 
evidence on human decision making, it is a powerful way to approximate 
intelligent behavior that has opened up a huge range of empirical applications.

Despite the common ancestry, there are clear differences in the practical reasons for doing DDC and IRL modeling.
In economics, DDC and other types of structural models (i.e. dynamic, stochastic
general equilibrium models in macroeconomics)  are used to provide counterfactual
simulations to assess the impact of a wide variety of change in economic
policies on both individual and colective outcomes, behavior, and welfare. 
Structural models overcome a limitation 
of reduced form models, which treat policy variables as a covariate and rely on historical variation in policies to make
counterfactual predictions. Structural models are able to produce counterfactual predictions
of the impact of novel policies that have never been enacted or even considered previously.
This is why structural models are increasingly used for policymaking because it is cheaper and faster to ``crash test'' new policy ideas
via counterfactual simulations in an ``offline'' environment rather
than trying out new policies ``online'' and learning by trial and error.

In the AI and ML literatures, IRL has proven useful as a way to supply 
a reward function used by Reinforcement Learning (RL) algorithms for offline training of intelligent algorithms in robotics and AI. 
There is a wide range of complex, hard to formalize tasks where the appropriate reward function
is unclear {\it a priori.\/} IRL is able to infer this reward function from limited human ``demonstration data.'' 
For example, \citet{christiano:2017} notes that
``For sophisticated reinforcement learning (RL) systems to interact usefully with real-world environments, we need to communicate complex goals to these systems. In this work, we explore goals defined in terms of (non-expert) human preferences between pairs of trajectory segments. 
We show that this approach can effectively solve complex RL tasks without access to the reward function, including Atari games and simulated robot locomotion, while providing feedback on less than 1\% of our agent's interactions with the environment. 
This reduces the cost of human oversight far enough that it can be practically applied to state-of-the-art RL systems. 
To demonstrate the flexibility of our approach, we show that we can successfully train complex novel behaviors with about an hour of human time. These behaviors and environments are considerably more complex than any which have been previously learned from human feedback.''

Another important difference in the IRL and DDC literatures
is that the latter is more focused on statistical inference of the 
estimated reward function and deriving approximate distributions that reflect how
estimates are affected by sampling errors. As a result, many  DDC applications involve estimation of 
smaller, more parsimonious models with the goal of using them to gain new insights
into human decision making and test behavioral hypotheses.  The IRL literature is more focused on 
uncovering a reward function that can be used to train large scale applications using RL.  
IRL has been used to uncover reward functions for concrete problems such as route choice
that can require estimation of thousands and sometimes even millions of parameters. 
The large number of parameters and complexity of such reward
functions makes it challenging to evaluate them from the perspective of statistical inference and insight
into human  behavior. However the  ``proof is in the pudding'' --- 
the highly sophisticated and intelligent behavior in robotic and AI applications that have been learned
after extensive RL training using the reward functions provided by IRL (which are in turn learned from more limited observations of human behavior).

The DDC and IRL literatures have both greatly benefitted from the use of Deep Neural Networks (DNNs) to 
approximate value functions in problems with large state spaces. Though it is still not
formally established, the use of DNNs and other function
approximation  methods do seem to break Bellman's ``curse of dimensionality'' (at least in a practical sense) 
allowing us to formulate and estimate increasingly large and more realistic models.

A challenging area for further work is to relax and test the strong rationality assumptions underlying
DDC and IRL models. There are two key assumptions: 1) rational expectations (i.e. the subjective
beliefs about the decision-dependent law of motion for state variables coincides with the true one), and 2) the 
ability to perfectly optimize given one's beliefs. If these
assumptions do not hold, rewards inferred from human behavior will be distorted by our attempt to ``rationalize'' 
behavior that may not be fully rational. The notion that human behavior is more likely to be only ``boundedly rational''
goes back to the work of \citet{Simon:1957}, and it is exemplified by the fact that chess algorithms trained using RL 
such as Alpha-Zero consistently beat the best human Grand Masters. 
RL was able to do this without any need for human data and IRL
 because 1) the reward function is known {\it a priori\/} (i.e. 1 if win, 0 if lose), and 
2) the algorithm as trained via self-play.

The ``model-free'' approach to inferring rewards without explicitly
specifying beliefs that is used in some recent work in IRL contrasts largely model-based 
estimation approach in the DDC literature. Though estimation using fewer assumptions
is generally an attribute, model-free estimation does entail the strong implicit assumption that
agents we are modeling have rational expectations.  To the extent that humans are boundedly rational,
we would like to understand whether suboptimal behavior is due to distorted or irrational
beliefs or the inability to perfectly optimize.  In simpler environments that are less strategically complex than
chess, humans may be more likely to perfectly optimize but if their beliefs are not rational, 
the resulting behavior will  still be suboptimal from the standpoint of a rational observer.

Fro example, \citet{tennis:2025} showed that elite tennis professionals use suboptimal serve strategies,
and demonstrated how their behavior can be explained by distorted subjective beliefs 
about their own strengths and weaknesses as servers as well as their opponents. 
This opens a path for DDC and IRL models to improve human performance using data and
model-free approaches that enables them to correct distorted subjective beliefs. However
there are many high stake strategic situations (e.g. a wartime scenario)  where the
relevant data may be too scarce, leaving any decision maker no choice but to rely on an implicit
``mental model'' of the situation.  In such situations neither DDC nor IRL methods may
prove useful because of their dependence on good models of beliefs. More work needs to be 
done on how it might possible to construct relatively
accurate mental models (i.e. beliefs about how different actions affect outcomes), 
in unique situations where relevant data is extremely limited.



We conclude that the limitations confronting DDC and IRL may not be computational
but rather inherent limitations in our ability to infer underlying rewards and beliefs 
from data on states and actions as we noted in our discussion of the identification problem in 
section 2.4. The solution to the identification problem requires us to impose  strong 
prior restrictions (assumptions) about rewards and
beliefs. If these assumptions are wrong, counterfactual predictions from these
models may not be accurate even though the models may provide a good fit to the demonstration data
used to estimate them. However practical experience suggests that even if our assumptions
may not be precisely correct and human decision makers may not be not fully rational,
the assumptions used to identify DDC and IRL models appear to be sufficiently good approximations to 
produce realistic counterfactual simulations and leverage limited human demonstration data to enable us to train 
artificially intelligent agents to perform well in a variety of tasks in a variety
of unfamiliar situations that previously were tasks only human beings could do well.
The rapid growth in the DDC and IRL literatures suggests that despite the limitations
and challenges we have noted, these methods have proved to have enormous practical value
and so we are likely to see continued rapid progress  in the future that will enable these methods  
to be even more widely used to transform economics, finance and  artificial intelligence.





